\documentclass[12pt]{article}

\usepackage{amsmath,amssymb,amsthm,enumerate}
\usepackage{setspace}
\usepackage{hyperref}
\usepackage{geometry}
\usepackage{booktabs}
\usepackage{array}
\usepackage{mathtools}
\geometry{margin=1.2in}

\usepackage[round,comma,authoryear]{natbib}

\onehalfspacing

\DeclareMathOperator*{\argmax}{arg\,max}
\DeclareMathOperator*{\argmin}{arg\,min}

\newtheorem{theorem}{Theorem}
\newtheorem{proposition}{Proposition}
\newtheorem{lemma}{Lemma}
\newtheorem{corollary}{Corollary}
\newtheorem{remark}{Remark}
\newtheorem{definition}{Definition}

\newcommand{\calU}{\mathcal{U}}
\newcommand{\calZ}{\mathcal{Z}}
\newcommand{\zbar}{\bar{z}}
\newcommand{\zlow}{z_{0}}
\newcommand{\R}{\mathbb{R}}

\title{An HJB Equation for the Mirrlees Problem:\\
       A Structural Parallel with the Ramsey Tax Problem}
\author{Thomas J.\ Sargent}
\date{\today}

\begin{document}
\maketitle

\begin{abstract}
  This note constructs a Hamilton--Jacobi--Bellman (HJB) equation for the
  Mirrlees optimal nonlinear income-tax problem by exploiting a structural
  parallel with the continuous-time Ramsey tax problem studied in Jiang,
  Sargent, and Wang (2025).  In the Ramsey problem, decisions are ordered by
  time $t \ge 0$ and a single implementability constraint attaches a constant
  Lagrange multiplier $\Phi_0$ to the planner's problem; the value function is
  linear in a scalar state variable $X_t$.  In the Mirrlees problem, decisions
  are ordered by the household's productivity type $z \in [z_0, \bar{z}]$ and a
  single resource constraint attaches a constant multiplier $\mu$; the
  appropriate HJB has the value function linear in the state variable
  $\calU(z)$---the utility assigned to type $z$---under quasilinear
  preferences.  In the general case, the HJB still has a canonical form but
  the value function is not linear in $\calU$.  The note draws a
  correspondence table between the two problems.
\end{abstract}

\tableofcontents
\newpage

%% ===========================================================================
\section{The Ramsey Problem with No Initial Term Debt}
\label{sec:ramsey}
%% ===========================================================================

\subsection{Environment}

Following Jiang, Sargent, and Wang (2025), hereafter JSW, consider a
continuous-time Lucas--Stokey economy.  A representative household and a
benevolent government operate on $t \in [0,\infty)$.  The resource constraint
is
\begin{equation}\label{eq:R_resource}
  C_t + G_t = N_t, \qquad t \ge 0,
\end{equation}
where $C_t$ is consumption, $N_t$ is labor supply, and $G_t \ge 0$ is an
exogenous government expenditure flow.  The household's lifetime utility is
\begin{equation}\label{eq:R_utility}
  \int_0^\infty e^{-\rho t}\, U(C_t, N_t)\, dt, \qquad \rho > 0.
\end{equation}
Throughout this section we impose the special case:
\begin{equation}\label{eq:R_nodebt}
  b_{0-,t} = 0 \quad \text{for all } t \ge 0,
\end{equation}
meaning that the government carries \emph{no} initial term debt.

\subsection{The implementability constraint}

Using the household's first-order conditions and the government's budget
constraint, JSW show that an allocation $\{C_t, N_t\}$ is a competitive
equilibrium allocation if and only if it satisfies the resource constraint
\eqref{eq:R_resource} and the \emph{implementability constraint}.  Under
\eqref{eq:R_nodebt} this constraint simplifies to
\begin{equation}\label{eq:R_impl}
  \int_0^\infty e^{-\rho t}\bigl[C_t\, U_{C,t} + N_t\, U_{N,t}\bigr]\, dt = 0.
\end{equation}
The Ramsey planner selects the best feasible allocation by maximizing
$\int_0^\infty e^{-\rho t} U(C_t, N_t)\, dt$ subject to
\eqref{eq:R_resource}--\eqref{eq:R_impl}.

\subsection{The cumulative-liability state variable}

Because $b_{0-,t} = 0$, the government's cumulative liability $\Pi_t$
(the net financial position the government would carry if it never restructured
its debt) reduces to
\begin{equation}\label{eq:R_Pi}
  \Pi_t = \int_0^t \frac{q_{0,s}}{q_{0,t}}\bigl(G_s - \tau_s N_s\bigr)\, ds,
  \qquad \Pi_0 = 0,
\end{equation}
where $q_{0,t} = e^{-\rho t} U_{C,t}/U_{C,0}$ is the zero-coupon bond price
and $\tau_t$ is the flat labor-income tax rate.  In differential form,
\begin{equation}\label{eq:R_dPi}
  d\Pi_t = r_t\,\Pi_t\, dt + (G_t - \tau_t N_t)\, dt,
\end{equation}
where $r_t = \rho + \dot{C}_t/C_t$ is the equilibrium short rate.

JSW define the scalar state variable
\begin{equation}\label{eq:R_Xt}
  X_t = \Pi_t\, U_{C,t},
\end{equation}
which encodes both the government's cumulative liability and the household's
marginal utility of consumption.  Differentiating and using \eqref{eq:R_dPi}
together with the household's first-order conditions, they obtain
\begin{equation}\label{eq:R_dXt}
  \dot{X}_t = \rho\, X_t - C_t\, U_{C,t} - N_t\, U_{N,t},
\end{equation}
where the $b_{0-,t}$ terms present in the general JSW formula vanish under
\eqref{eq:R_nodebt}.  Note that the dynamics are \emph{linear} in $X_t$ (the
$\rho X_t$ term).

\medskip

\noindent\textbf{Initial condition.}  Because $\Pi_0 = 0$, the state starts at
$X_0 = 0$.

\subsection{The HJB equation}

Define the value function
\begin{equation}\label{eq:R_V}
  V(X_t, t) = \max_{\{C_s;\, s \ge t\}} \int_t^\infty e^{-\rho(s-t)} U(C_s, C_s+G_s)\, ds,
\end{equation}
subject to the law of motion \eqref{eq:R_dXt}.  By the principle of dynamic
programming, $V$ satisfies the HJB equation
\begin{equation}\label{eq:R_HJB}
  \rho\, V = V_t + \max_{C_t}\Bigl\{U(C_t, C_t+G_t)
  + V_X\,\bigl(\rho X_t - C_t U_{C,t} - N_t U_{N,t}\bigr)\Bigr\}.
\end{equation}

\subsection{The linear value-function guess and its verification}

\noindent\textbf{Guess.}  Try the additively separable form
\begin{equation}\label{eq:R_Vguess}
  V(X_t, t;\, \Phi_0) = -\Phi_0\, X_t + H(t;\, \Phi_0),
\end{equation}
where $\Phi_0 \ge 0$ is a scalar and $H$ is a function of time and $\Phi_0$.
Then $V_X = -\Phi_0$ (constant in both $X_t$ and $t$).

\medskip\noindent\textbf{Verification.}  Substituting \eqref{eq:R_Vguess}
into \eqref{eq:R_HJB}, the $-\Phi_0\, \rho X_t$ terms on both sides cancel,
and the HJB reduces to the following ODE for $H$:
\begin{equation}\label{eq:R_Hode}
  \rho\, H(t;\, \Phi_0) = H'(t;\, \Phi_0)
  + \max_{C_t}\Bigl\{U(C_t, C_t+G_t)
  + \Phi_0\bigl(C_t U_{C,t} + N_t U_{N,t}\bigr)\Bigr\}.
\end{equation}
The maximization over $C_t$ on the right side is independent of the state
$X_t$.  The optimal consumption satisfies
\begin{equation}\label{eq:R_Copt}
  C(\Phi_0;\, G_t) = \argmax_{C} \Bigl\{U(C, C+G_t)
  + \Phi_0\bigl(C\, U_C + (C+G_t) U_N\bigr)\Bigr\}.
\end{equation}

\subsection{Determining the multiplier $\Phi_0$}

Substituting the optimal plan $C^* = C(\Phi_0; G_t)$ and $N^* = C^* + G_t$ into
the implementability constraint \eqref{eq:R_impl} and imposing $X_0 = 0$,
the Ramsey multiplier $\Phi_0^*$ is the non-negative root of
\begin{equation}\label{eq:R_IC}
  I_0(\Phi_0) \equiv
  \int_0^\infty e^{-\rho t}\bigl[C(\Phi_0;\, G_t)\, U_C + N(\Phi_0;\, G_t)\, U_N\bigr]\, dt
  = 0,
\end{equation}
where $N(\Phi_0; G_t) = C(\Phi_0; G_t) + G_t$.

\medskip
\begin{remark}[Economic interpretation of $\Phi_0$]
The multiplier $\Phi_0$ is the shadow price of the single implementability
constraint and is \emph{constant across all dates}.  This constancy reflects
the exponential discount structure in the Ramsey problem: the $\rho X_t$ term
in the law of motion \eqref{eq:R_dXt} exactly cancels the discounting on the
left side of the HJB, keeping the shadow price $-V_X = \Phi_0$ time-invariant.
\end{remark}

\medskip
\begin{remark}[Constancy of the optimal tax rate when $G_t = G$]
When government spending is constant ($G_t = G$ for all $t$), the
right-hand side of \eqref{eq:R_Copt} does not depend on $t$, so
$C^* = C(\Phi_0^*; G)$ is constant.  The implementability constraint
\eqref{eq:R_IC} then collapses to the static condition
$C^* U_{C}^* + N^* U_{N}^* = 0$, and the Ramsey tax rate is constant at
$\tau^* = G/N^*$.  This echoes Barro's (1979) tax-smoothing result.
\end{remark}

%% ===========================================================================
\section{The Mirrlees Problem: Sequential Structure in Productivity}
\label{sec:mirrlees}
%% ===========================================================================

\subsection{Environment}

Following the lecture notes on optimal nonlinear taxation, there is a continuum
of households indexed by productivity $z \in \calZ = [\zlow, \zbar]$ with
density $f(z) > 0$.  A type-$z$ household produces output
\begin{equation}\label{eq:M_output}
  y(z) = z\, n(z),
\end{equation}
where $n(z)$ is labor and $z$ is the marginal product.  Utility over
consumption $c$ and labor $n$ is $U(c, n)$.  The government observes only
income $y$, not the decomposition $(z, n)$.

\subsection{The mechanism design formulation}

By the taxation principle, the government can equivalently optimize over
incentive-compatible (IC) allocations $\{c(z), y(z)\}$.  Defining household
utility
\begin{equation}\label{eq:M_utility}
  \calU(z) = U\!\left(c(z),\, \frac{y(z)}{z}\right),
\end{equation}
the local IC condition (envelope theorem applied to the announcement problem)
is
\begin{equation}\label{eq:M_IC}
  \frac{d\calU}{dz} = -U_n\!\left(c(z),\,\frac{y(z)}{z}\right)\frac{y(z)}{z^2}
  \equiv \phi(z),
\end{equation}
where $\phi(z) > 0$ because $U_n < 0$.  Condition \eqref{eq:M_IC} is the
\emph{envelope ODE}: it plays the role of the law of motion for the state
variable when the problem is viewed sequentially in $z$.

Consumption is the implicit function of state and control:
\begin{equation}\label{eq:M_c}
  c(z) = C\!\bigl(\calU(z),\, y(z)\bigr),
  \qquad
  \text{where } \calU = U(C(\calU, y),\, y/z).
\end{equation}
Its partial derivatives are $C_\calU = 1/U_c$ and $C_y = -U_n/(z U_c)$.

\subsection{The planner's problem}

The government maximizes social welfare subject to the IC law of motion
\eqref{eq:M_IC} and a resource constraint:
\begin{alignat}{2}
  & \max_{y(\cdot),\, \calU(\zlow)} &
  & \int_{\zlow}^{\zbar} \psi(z)\,\calU(z)\,f(z)\,dz
  \label{eq:M_obj} \\
  & \text{s.t.} &
  & \frac{d\calU}{dz} = -U_n\!\bigl(C(\calU,y),\, y/z\bigr)\frac{y}{z^2},
  \label{eq:M_IC2} \\
  & &
  & \int_{\zlow}^{\zbar}\bigl[y(z)-C(\calU(z), y(z))\bigr]f(z)\,dz \ge B.
  \label{eq:M_RC}
\end{alignat}
The control is the output schedule $y(\cdot)$; the state is $\calU(z)$.  The
initial utility level $\calU(\zlow)$ is a free endpoint to be optimized.

The single \emph{global} constraint \eqref{eq:M_RC} plays the same structural
role as the implementability constraint \eqref{eq:R_impl} in the Ramsey
problem.  Attaching a non-negative Lagrange multiplier $\mu$ to \eqref{eq:M_RC},
the augmented objective is
\begin{equation}\label{eq:M_augmented}
  \max_{y(\cdot),\, \calU(\zlow)}
  \int_{\zlow}^{\zbar}
  \bigl[\psi(z)\,\calU(z) + \mu\bigl(y(z) - C(\calU(z), y(z))\bigr)\bigr]
  f(z)\,dz - \mu B.
\end{equation}

%% ===========================================================================
\section{The Mirrlees HJB Equation}
\label{sec:mirrlees_hjb}
%% ===========================================================================

\subsection{The sequential value function}

Viewing productivity $z$ as the ``index'' along which the planner proceeds
(just as $t$ indexes the Ramsey planner's decisions), define the
\emph{continuation value} for the planner who has already assigned allocations
to types below $z$ and now faces types $[\zlow, \ldots] \to [z, \zbar]$ with
inherited state $\calU(z) = u$:
\begin{align}
  V(u, z;\, \mu) = \max_{y(\cdot)}
  \int_z^{\zbar}
  \bigl[\psi(\zeta)\,\calU(\zeta) + \mu\bigl(y(\zeta) - C(\calU(\zeta), y(\zeta))\bigr)\bigr]
  f(\zeta)\,d\zeta,
  \label{eq:M_V}
\end{align}
subject to \eqref{eq:M_IC2} and $\calU(z) = u$.

\medskip\noindent\textbf{Commitment through promised values.}
The state variable $u = \calU(z)$ is the utility that has already been
\emph{promised} to type $z$ by the assignments made to all lower types
$\zeta < z$.  A continuation planner starting at $z$ is therefore bound to
deliver exactly $u$ to type $z$: it inherits a commitment it cannot revise.
In this sense the sequential-in-$z$ formulation embeds full commitment by the
original mechanism designer.  Types already processed ($\zeta < z$) have
received their allocations and cannot be renegotiated; the promised value $u$
is the only ``footprint'' those past decisions leave on the continuation
problem.  This is the Mirrlees analogue of the Ramsey planner's commitment to
an initial household budget constraint: just as the state variable $X_t =
\Pi_t U_{C,t}$ in the Ramsey HJB records the cumulative fiscal obligation
that the Ramsey planner cannot escape, the promised utility $u$ in the
Mirrlees HJB records the incentive obligation to type $z$ that the Mirrlees
planner cannot escape.

\subsection{The HJB equation}

By the principle of dynamic programming, $V$ satisfies
\begin{equation}\label{eq:M_HJB}
  -\frac{\partial V}{\partial z} =
  \max_{y}\Bigl\{
  \bigl[\psi(z)\, u + \mu\bigl(y - C(u, y)\bigr)\bigr]f(z)
  + \frac{\partial V}{\partial u}\cdot
  \Bigl(-U_n\!\bigl(C(u,y),\,y/z\bigr)\frac{y}{z^2}\Bigr)
  \Bigr\}.
\end{equation}
This is the \emph{Mirrlees HJB equation}.  Writing $\phi(u, y, z) \equiv
-U_n(C(u,y), y/z)\, y/z^2 > 0$, \eqref{eq:M_HJB} can be compactly stated as
\begin{equation}\label{eq:M_HJBcompact}
  -V_z = \max_y \Bigl\{
  [\psi(z)\, u + \mu(y - C(u,y))] f(z)
  + V_u\cdot \phi(u, y, z)
  \Bigr\}.
\end{equation}

\noindent\textbf{Boundary condition.}  At the top of the type distribution
there is no continuation value, so
\begin{equation}\label{eq:M_HJBBC}
  V(u, \zbar;\, \mu) = 0 \quad \text{for all } u.
\end{equation}

\subsection{The linear value-function guess}

\noindent\textbf{Guess.}  Motivated by the Ramsey analogy, try
\begin{equation}\label{eq:M_Vguess}
  V(u, z;\, \mu) = \lambda(z;\, \mu)\cdot u + K(z;\, \mu),
\end{equation}
where $\lambda$ is a function of $z$ and $K$ is a separate function of $z$ and
$\mu$.  Then $V_u = \lambda(z)$.

\medskip\noindent\textbf{When does the guess work?}  The linear guess is valid
whenever the maximand in \eqref{eq:M_HJBcompact} is \emph{affine in $u$}.
This requires:
\begin{enumerate}[(i)]
  \item The flow payoff $[\psi(z) u + \mu(y - C(u,y))] f(z)$ to be affine in $u$, which holds when $C(u, y)$ is affine in $u$.
  \item The state dynamics $\phi(u, y, z)$ to be independent of $u$ (or affine in $u$), which holds when $U_n(C(u,y), y/z)$ is independent of $u$.
\end{enumerate}
Both conditions are satisfied under \emph{quasilinear preferences}, $U(c,n) =
c - v(n)$, since then $C(u, y) = u + v(y/z)$, which is linear in $u$ with
slope $1$, and $U_n = -v'(y/z)$ is independent of $c$ and hence of $u$.

\medskip\noindent\textbf{Verification under quasilinear preferences.}  With
$U = c - v(n)$:
\begin{align*}
  C(u, y) &= u + v(y/z), \\
  \phi(u, y, z) &= v'(y/z)\, y/z^2.
\end{align*}
Substituting the linear guess \eqref{eq:M_Vguess} into \eqref{eq:M_HJBcompact}:
\begin{multline*}
  -\lambda'(z)\, u - K'(z) \\
  = \max_y \Bigl\{
  \bigl[\psi(z) - \mu\bigr] u\, f(z)
  + \mu\bigl[y - v(y/z)\bigr] f(z)
  + \lambda(z)\, v'(y/z)\, y/z^2
  \Bigr\}.
\end{multline*}
Separating terms in $u$ from terms independent of $u$:
\begin{align}
  \lambda'(z) &= \bigl[\mu - \psi(z)\bigr] f(z),
  \label{eq:M_costate} \\[4pt]
  K'(z) &= -\max_{y}\Bigl\{\mu\bigl[y - v(y/z)\bigr] f(z)
             + \lambda(z)\, v'(y/z)\, y/z^2\Bigr\}.
  \label{eq:M_Kode}
\end{align}
Equation \eqref{eq:M_costate} is exactly the costate equation derived from the
Hamiltonian approach in the lecture notes.

\subsection{Optimality condition for the output schedule}

The first-order condition for the maximization problem in \eqref{eq:M_Kode}
(in the quasilinear case) is
\begin{equation}\label{eq:M_ystar}
  0 = \mu\bigl[1 - v'(y/z)/z\bigr] f(z)
  + \frac{\lambda(z)}{z^2}\bigl[v'(y/z) + n\, v''(y/z)\bigr],
\end{equation}
where $n = y/z$.  Rearranging using $v'(n) = z(1-T'(y))$:
\begin{equation}\label{eq:M_taxformula}
  \frac{T'(y(z))}{1-T'(y(z))}
  = \Bigl[1 + \frac{n\, v''(n)}{v'(n)}\Bigr]\cdot
  \frac{\lambda(z)}{z\, f(z)\, \mu}.
\end{equation}
This is the Diamond (1998) ABC formula.

\subsection{Boundary conditions and determination of $\mu$}

\medskip\noindent\textbf{At the top.}  The boundary condition \eqref{eq:M_HJBBC}
with the guess \eqref{eq:M_Vguess} gives
\begin{equation}\label{eq:M_bctop}
  \lambda(\zbar) = 0, \qquad K(\zbar) = 0.
\end{equation}

\medskip\noindent\textbf{At the bottom.}  The initial utility $\calU(\zlow)$
is a free endpoint (the planner can choose the utility of the lowest type).
The transversality condition at the free endpoint is $V_u(u, \zlow) = 0$, which
by the guess \eqref{eq:M_Vguess} gives
\begin{equation}\label{eq:M_bcbot}
  \lambda(\zlow) = 0.
\end{equation}

Together, \eqref{eq:M_bctop}--\eqref{eq:M_bcbot} are the two-point boundary
conditions on the costate ODE \eqref{eq:M_costate}.  They are exactly the
transversality conditions $\lambda(\zlow) = \lambda(\zbar) = 0$ derived in the
lecture notes.  In the quasilinear case, integrating \eqref{eq:M_costate} from
$z$ to $\zbar$ using $\lambda(\zbar) = 0$ gives
\begin{equation}\label{eq:M_lambdasol}
  \lambda(z) = \int_z^{\zbar} \bigl[\psi(\zeta) - \mu\bigr] f(\zeta)\, d\zeta,
\end{equation}
and the lower boundary condition $\lambda(\zlow) = 0$ pins down $\mu = 1$
(the marginal utility of public funds equals one under quasilinear
preferences).

\medskip\noindent\textbf{Determining $\mu$ in general.}  For general
preferences, the multiplier $\mu$ is determined by the resource constraint
\eqref{eq:M_RC}:
\begin{equation}\label{eq:M_mudet}
  \int_{\zlow}^{\zbar}\bigl[y^\ast(z) - C(\calU^\ast(z), y^\ast(z))\bigr] f(z)\, dz = B,
\end{equation}
which plays the role of equation \eqref{eq:R_IC} in the Ramsey problem.

%% ===========================================================================
\section{Structural Correspondence between the Two Problems}
\label{sec:correspondence}
%% ===========================================================================

Table~\ref{tab:correspondence} displays the element-by-element structural
parallel between the Ramsey problem (no initial term debt) and the Mirrlees
problem.

\begin{table}[h!]
\centering
\caption{Structural Correspondence: Ramsey vs.\ Mirrlees}
\label{tab:correspondence}
\renewcommand{\arraystretch}{1.45}
\begin{tabular}{@{} p{6.2cm} p{6.2cm} @{}}
\toprule
\textbf{Ramsey problem (no initial debt)} & \textbf{Mirrlees problem} \\
\midrule
Sequential index: time $t \in [0,\infty)$ &
Sequential index: productivity $z \in [\zlow, \zbar]$ \\

State variable: $X_t = \Pi_t\, U_{C,t}$ &
State variable: $\calU(z)$ (household utility) \\

Control: consumption $C_t$ &
Control: output schedule $y(z)$ \\

Law of motion:
$\dot{X}_t = \rho X_t - C_t U_{C,t} - N_t U_{N,t}$ &
Envelope ODE:
$\calU'(z) = -U_n(c, y/z)\, y/z^2$ \\

Initial condition: $X_0 = 0$ &
Free endpoint: $\calU(\zlow)$ optimized; $\lambda(\zlow)=0$ \\

Terminal condition: transversality as $t\to\infty$ &
Terminal condition: $\lambda(\zbar)=0$ \\

Objective:\newline $\displaystyle\int_0^\infty e^{-\rho t} U(C_t, N_t)\, dt$ &
Objective:\newline $\displaystyle\int_{\zlow}^{\zbar} \psi(z)\, \calU(z)\, f(z)\, dz$ \\

Single global constraint: implementability ($=0$) &
Single global constraint: resource ($=B$) \\

Global-constraint multiplier $\Phi_0$ &
Global-constraint multiplier $\mu$ \\

Value function: $V = -\Phi_0 X_t + H(t;\, \Phi_0)$ &
Value function: $V = \lambda(z)\, \calU + K(z;\, \mu)$\newline
(linear in $\calU$ under quasilinear preferences) \\

HJB shadow price: $V_X = -\Phi_0$ \emph{(constant)} &
HJB shadow price: $V_\calU = \lambda(z)$ \emph{(varies with $z$)} \\

Costate dynamics:
$-\Phi_0 = \text{const}$ &
Costate ODE:
$\lambda'(z) = [\mu - \psi(z)] f(z)$ \\

Optimal control: $C(\Phi_0;\, G_t) = \displaystyle\argmax_C [\ldots]$ &
Optimal control: $y(z) = \argmax_y [\ldots]$ (tax formula) \\

Determines $\Phi_0$: $I_0(\Phi_0) = 0$ &
Determines $\mu$: resource constraint $= B$ \\
\bottomrule
\end{tabular}
\end{table}

\medskip
A key structural difference deserves emphasis:

\begin{remark}[Constant vs.\ varying shadow price]
In the Ramsey problem, the shadow price $V_X = -\Phi_0$ is the same
constant at every date $t$.  This constancy arises because the state dynamics
\eqref{eq:R_dXt} contain a $\rho X_t$ term that cancels the discounting
$\rho V$ on the left of the HJB \eqref{eq:R_HJB}, leaving the costate
equation trivially satisfied with $\Phi_0 = \text{const}$.

In the Mirrlees problem, the shadow price $\lambda(z) = V_\calU$ evolves with
$z$ according to the costate ODE \eqref{eq:M_costate}.  The reason is that
the envelope ODE \eqref{eq:M_IC} contains no term proportional to $\calU$
(analogous to no ``$z$-discounting''), and the welfare weight $\psi(z)$ varies
with $z$.  The costate $\lambda(z)$ accumulates the gap between the social
welfare weight $\psi$ and the social cost of public funds $\mu$, measuring the
``redistribution pressure'' emanating from types above $z$.
\end{remark}

\begin{remark}[Quasilinear preferences as the exact Ramsey analogue]
In the Ramsey problem, the linear value-function guess $V = -\Phi_0 X_t +
H(t)$ works for \emph{all} admissible preferences.  This is because the
dynamics \eqref{eq:R_dXt} are linear in $X_t$ for any preferences, and the
optimization over $C_t$ in the HJB is independent of $X_t$.

In the Mirrlees problem, the analogous linear guess $V = \lambda(z) \calU +
K(z)$ works only under \emph{quasilinear preferences} $U = c - v(n)$, because
only then is the implicit function $C(\calU, y) = \calU + v(y/z)$ linear in
$\calU$, making the HJB affine in the state $\calU$.  For general preferences,
$C(\calU, y)$ is nonlinear in $\calU$, and the value function $V(u,z)$ is
nonlinear in $u$.  The HJB \eqref{eq:M_HJBcompact} remains the correct
dynamic-programming equation in both cases; it is only the validity of the
\emph{linear} guess that is special to quasilinear preferences.
\end{remark}

%% ===========================================================================
\section{The Mirrlees HJB in the General Case}
\label{sec:general}
%% ===========================================================================

For general preferences, the HJB \eqref{eq:M_HJBcompact} still governs the
value function, but $V$ is a nonlinear function of $u$.  The optimality
conditions derived from the Hamiltonian approach in the lecture notes---the
$H_y = 0$ condition and the costate equation---are exactly the first-order
conditions of the HJB: they arise from differentiating the maximand in
\eqref{eq:M_HJBcompact} with respect to $y$ and equating the $z$-derivative
of $V_\calU$ to minus the $\calU$-gradient of the maximand.

Specifically, denote $\lambda(z) \equiv V_u(u^*(z), z)$ along the optimal
path, where $u^*(z) = \calU^*(z)$ is the optimal utility assignment.  Then:

\begin{enumerate}[(i)]

\item \textbf{Optimality in $y$} (from $\partial/\partial y$ of the maximand):
\begin{equation}\label{eq:M_Hy_general}
  \mu \, [1 - C_y] f(z)
  - \lambda(z)\Bigl[U_{nc}\, C_y + U_{nn}/z\Bigr]\frac{y}{z^2}
  - \lambda(z)\frac{U_n}{z^2} = 0.
\end{equation}

\item \textbf{Costate equation} (from $-\partial/\partial u$ of the maximand):
\begin{equation}\label{eq:M_costate_general}
  \frac{d\lambda}{dz} = \mu\, C_\calU\, f(z) - \psi(z)\, f(z)
  + \lambda(z)\, U_{nc}\,\frac{y}{z^2}\, C_\calU.
\end{equation}

\end{enumerate}
These are identical to equations (5.2) and (5.3) of the lecture notes.  The HJB
perspective makes clear that $\lambda(z)$ is not an auxiliary variable of the
Hamiltonian approach but is the \emph{slope of the value function}
$\lambda(z) = V_u(\calU^*(z), z)$ evaluated along the optimal path.

%% ===========================================================================
\section{Summary and Interpretation}
\label{sec:summary}
%% ===========================================================================

Both the Ramsey and Mirrlees planning problems share the following structural
architecture:

\begin{enumerate}[(1)]

\item \textbf{A sequential ordering principle.}  In Ramsey, decisions are made
along the time axis $t$.  In Mirrlees, the social planner processes
households from $z = \zlow$ upward through higher productivities.  In both
cases, the planner integrates a \emph{flow payoff} and a \emph{law of motion}
for a one-dimensional state variable.

\item \textbf{A single global constraint that attaches a constant multiplier.}
In Ramsey, the implementability constraint \eqref{eq:R_impl} attaches the
constant $\Phi_0$.  In Mirrlees, the resource constraint \eqref{eq:M_RC}
attaches the constant $\mu$.  In both cases, this multiplier folds the global
constraint into the flow payoff and turns the problem into a standard
(undiscounted or discounted) calculus-of-variations problem.

\item \textbf{A Hamilton--Jacobi--Bellman equation.}  Once the global
constraint is Lagrangianized, the problem admits a value function that
satisfies an HJB.  The shadow price of the state variable ($V_X = -\Phi_0$ in
Ramsey; $V_\calU = \lambda(z)$ in Mirrlees) is the costate variable of the
underlying optimal control problem.

\item \textbf{A determining equation for the global multiplier.}  The multiplier
is found by imposing the global constraint: $I_0(\Phi_0) = 0$ in Ramsey, and
the resource constraint \eqref{eq:M_mudet} in Mirrlees.

\end{enumerate}

\noindent The deepest difference is that in Ramsey the costate is constant
(because exponential discounting with rate $\rho$ exactly offsets the $\rho
X_t$ drift in the state), while in Mirrlees the costate $\lambda(z)$ varies,
accumulating the difference between the planner's welfare weights and the
social cost of funds.  Under quasilinear preferences the linear value-function
guess works in both problems, and the two HJBs are fully isomorphic.

%% ===========================================================================
\section{Extension: Cyclical Government Expenditures in the Ramsey Problem}
\label{sec:cyclical_ramsey}
%% ===========================================================================

\subsection{Setup}

This section studies the Ramsey HJB under cyclical government expenditures,
mirroring Section~5.3 of JSW.  Maintain the no-initial-debt specialization
$b_{0-,t}=0$ and adopt the utility function $U(C,N)=\ln C - N$ used in JSW's
quantitative examples.  Government spending follows
\begin{equation}\label{eq:cyc_G}
  G_t = G_0 + \theta_G \cos(\omega t), \qquad G_0,\,\theta_G>0,\; \omega>0,
\end{equation}
so that fiscal needs oscillate around the mean $G_0$ with amplitude $\theta_G$
and angular frequency $\omega$.

\subsection{Constant optimal tax from the HJB}

Under $U = \ln C - N$ we have $U_C = 1/C$ and $U_N = -1$.  With
$b_{0-,t}=0$, the maximization problem \eqref{eq:R_Copt} in the HJB becomes
\begin{equation}\label{eq:cyc_HJBmax}
  C(\Phi_0;\, G_t)
  = \argmax_{C>0}\Bigl\{
  \ln C - (C+G_t)
  + \Phi_0\bigl[C\cdot\tfrac{1}{C} + (C+G_t)(-1)\bigr]
  \Bigr\}
  = \argmax_{C>0}\Bigl\{
  \ln C - (1+\Phi_0)(C+G_t)
  \Bigr\}.
\end{equation}
The first-order condition $1/C = 1+\Phi_0$ yields
\begin{equation}\label{eq:cyc_Cstar}
  C^* = \frac{1}{1+\Phi_0^*},
  \qquad
  \tau^* = 1 - C^* = \frac{\Phi_0^*}{1+\Phi_0^*}.
\end{equation}
Crucially, $C^*$ and $\tau^*$ are \emph{independent of $G_t$}: the cyclical
component of fiscal needs does not appear in the static optimality condition
for consumption.  The Ramsey planner therefore imposes a \emph{constant} tax
rate on every date, even though government expenditures oscillate.  This is
the HJB incarnation of Barro's (1979) tax-smoothing result.

\medskip
\begin{remark}[Why $G_t$ drops out]
The HJB maximand contains $G_t$ twice: once in the felicity term $-(C+G_t)$ and once
in the incentive term $-(1+\Phi_0)G_t$.  Both enter additively as constants
in the maximization over $C$.  The optimum $C^*$ is therefore pinned by the
ratio of the marginal utility of consumption to the Lagrange multiplier, with
no role for $G_t$.
\end{remark}

\subsection{Determining $\Phi_0^*$ from the implementability constraint}

The implementability constraint \eqref{eq:R_IC} with $U_C = 1/C^*$ and
$U_N = -1$ reduces to
\begin{equation}\label{eq:cyc_IC}
  \int_0^\infty e^{-\rho t}\bigl[1 - (C^*+G_t)\bigr]\, dt = 0,
\end{equation}
which gives
\begin{equation}\label{eq:cyc_PV}
  \frac{1-C^*}{\rho}
  = \int_0^\infty e^{-\rho t} G_t\, dt
  = \frac{G_0}{\rho}
  + \theta_G \frac{\rho}{\rho^2+\omega^2}.
\end{equation}
Hence
\begin{equation}\label{eq:cyc_tau}
  1-C^* = G_0 + \theta_G\frac{\rho^2}{\rho^2+\omega^2}.
\end{equation}
Three observations follow.  First, the equilibrium tax rate depends only on
the discounted present value of government expenditures, not on the timing of
the oscillations.  Second, the contribution of the cyclical component is
attenuated by the factor $\rho^2/(\rho^2+\omega^2) \in (0,1)$: high-frequency
fiscal cycles ($\omega \gg \rho$) contribute negligibly to the tax rate
because they nearly cancel in present value.  Third, as $\theta_G\to 0$ the
formula collapses to the constant-spending result $1-C^* = G_0$.

\subsection{The cyclical state variable $X_t^*$}

Even though $C^*$ and $\tau^*$ are constant, the state variable
$X_t^* = \Pi_t^* U_{C,t}^* = \Pi_t^*/C^*$ inherits the cyclicality of $G_t$.
Substituting $C^* = \text{const}$, $N_t^* = C^*+G_t$ into the law of motion
\eqref{eq:R_dXt}:
\begin{equation}\label{eq:cyc_dX}
  \dot X_t = \rho X_t - 1 + C^* + G_t
  = \rho X_t - \frac{\theta_G\rho^2}{\rho^2+\omega^2}
    + \theta_G\cos(\omega t),
\end{equation}
where we used $1-C^* = G_0 + \theta_G\rho^2/(\rho^2+\omega^2)$ to absorb
$G_0$.  The ODE \eqref{eq:cyc_dX} is a linear constant-coefficient equation
with a sinusoidal forcing term and a constant term that exactly offsets the
particular solution's mean.  Solving with $X_0 = 0$:
\begin{equation}\label{eq:cyc_Xsol}
  X_t^*
  = \frac{\theta_G}{\rho^2+\omega^2}
    \bigl[\rho\bigl(1-\cos(\omega t)\bigr) + \omega\sin(\omega t)\bigr].
\end{equation}
This solution is bounded, oscillates with period $2\pi/\omega$, and starts
and returns to values near zero.  For $\omega \gg \rho$ the amplitude of
oscillation is approximately $\theta_G/\omega$, so high-frequency spending
shocks produce only small oscillations in the state variable.

\medskip
\begin{remark}[$X_t^*$ as a fiscal buffer]
The state variable $X_t^*$ is the product of the government's cumulative
liability $\Pi_t^*$ and the household's marginal utility $1/C^*$.  Under
constant $C^*$, variations in $X_t^*$ map one-for-one into variations in
$\Pi_t^*$.  Equation~\eqref{eq:cyc_Xsol} shows that $\Pi_t^*$ absorbs all
cyclical fiscal fluctuations: when $G_t$ is high the government runs a deficit
(the primary surplus $S_t^* = \tau^* N_t^* - G_t$ falls), pushing $\Pi_t^*$
up; when $G_t$ returns to its mean, the surplus rises and $\Pi_t^*$ falls back.
The tax rate is constant throughout because it is the \emph{state variable},
not the tax rate, that adjusts.
\end{remark}

\subsection{Structure of $H(t;\, \Phi_0^*)$}

The time-varying component of the value function, $H(t;\,\Phi_0^*)$, satisfies
the ODE \eqref{eq:R_Hode} with $C_t = C^*$ fixed. Evaluating the maximand:
\begin{equation}\label{eq:cyc_Hode}
  \rho H(t;\,\Phi_0^*)
  = H'(t;\,\Phi_0^*)
  + \ln C^* - (C^*+G_t) + \Phi_0^*\bigl[1 - (C^*+G_t)\bigr].
\end{equation}
Because $G_t = G_0 + \theta_G\cos(\omega t)$, the right-hand side oscillates
in $t$.  The solution $H(t;\,\Phi_0^*)$ is therefore a sum of a constant
steady-state component (driven by $G_0$) and an oscillatory transient (driven
by $\theta_G\cos(\omega t)$):
\begin{equation}
  H(t;\,\Phi_0^*) = H_\infty + A_H\cos(\omega t) + B_H\sin(\omega t),
\end{equation}
where $H_\infty = [\ln C^* - (1+\Phi_0^*) C^*]/\rho$ is the constant
component and $(A_H, B_H)$ are determined by matching $\omega$-frequency
terms.  Although $H$ oscillates, the optimal $C$ does not: the oscillations
in $H$ represent the time-varying cost of financing cyclical government
expenditures, holding the tax plan fixed.  The value function
$V(X_t^*, t) = -\Phi_0^* X_t^* + H(t;\Phi_0^*)$ therefore oscillates both
through the state $X_t^*$ and directly through $H(t)$, the two oscillations
together reflecting the complete welfare cost of the cyclical fiscal program.

\medskip
\begin{remark}[Comparison with the constant-$G$ case]
When $\theta_G = 0$ the problem is fully stationary: $C^* = 1/(1+\Phi_0^*)$,
$\tau^* = G_0/N^* = G_0/(C^*+G_0)$ is constant, $r^* = \rho$ is constant,
$X_t^* = 0$ for all $t$, and $H(t;\Phi_0^*) = H_\infty$ is constant.  The
HJB collapses to a static saddle-point problem.  Introducing $\theta_G > 0$
breaks stationarity of $H$ and $X^*$ but leaves the control rule intact: the
optimal plan ``reacts'' to fiscal cycles entirely through the state variable,
not through the tax rate.
\end{remark}

%% ===========================================================================
\section{Extension: The Mirrlees HJB with Pareto-Distributed Productivity}
\label{sec:pareto}
%% ===========================================================================

\subsection{Setup and motivation}

The canonical Mirrlees analysis with bounded support $\calZ = [\zlow, \zbar]$
yields a \emph{zero top marginal tax rate}: $T'(y(\zbar)) = 0$.  This result
follows directly from the boundary condition $\lambda(\zbar) = 0$ together
with the tax formula \eqref{eq:M_taxformula}: at the top, $(1-F(\zbar)) = 0$
so the costate is zero and the distortion vanishes.  A large empirical
literature documents that the upper tail of the earnings distribution is
well-approximated by a Pareto distribution, not a distribution with bounded
support.  Diamond and Saez (2011) argue that the zero-top-rate result is an
artifact of bounded support and derive a formula for the optimal \emph{limiting}
marginal tax rate under an unbounded Pareto tail.

This section extends the HJB framework to unbounded support
$\calZ = [\zlow, \infty)$ and derives the Diamond--Saez top-rate formula
directly from the asymptotic behavior of the Mirrlees HJB.

\subsection{The Pareto distribution}

Let productivity $z$ follow a Pareto distribution with shape $\alpha > 1$
and lower bound $\zlow$:
\begin{equation}\label{eq:pareto_f}
  f(z) = \alpha\,\zlow^\alpha\, z^{-\alpha-1}, \qquad
  1 - F(z) = \left(\frac{\zlow}{z}\right)^\alpha, \qquad z \ge \zlow.
\end{equation}
The inverse hazard rate is
\begin{equation}\label{eq:pareto_haz}
  \frac{1-F(z)}{z\,f(z)} = \frac{1}{\alpha}\quad \text{for all } z \ge \zlow,
\end{equation}
a constant.  Under the optimal tax schedule, earnings $y(z)$ are also Pareto
in the tail.  With quasilinear preferences $U = c - v(n)$ and constant
elasticity $v(n) = n^{1+1/e}/(1+1/e)$, the earnings Pareto parameter is
\begin{equation}\label{eq:pareto_y}
  \alpha_y = \frac{\alpha\, e}{1+e},
\end{equation}
where $e = 1/\gamma$ is the compensated labor-supply (or taxable-income)
elasticity.

\subsection{Transversality replaces the upper boundary condition}

With $\zbar = \infty$, there is no upper type to impose $\lambda(\infty) = 0$
as a point condition.  Instead, optimality requires the transversality
condition
\begin{equation}\label{eq:pareto_TV}
  \lim_{z\to\infty} V(u,\, z;\, \mu) = 0 \quad \text{for all } u,
\end{equation}
meaning no continuation value remains once all types have been assigned
allocations.  Under the linear guess $V = \lambda(z)\,u + K(z;\mu)$, this
decomposes into
\begin{equation}\label{eq:pareto_BC}
  \lambda(\infty) = 0, \qquad K(\infty) = 0.
\end{equation}
These are exactly the conditions that ensure the costate integral
\eqref{eq:M_lambdasol} converges.  For the Pareto distribution, the integrand
$[\psi(\zeta)-\mu]f(\zeta)$ decays as $\zeta^{-\alpha-1}$, fast enough for
$\alpha > 1$, so the integral converges and $\lambda(\infty) = 0$ is
automatically satisfied as a transversality condition rather than imposed as
a point constraint.

\subsection{Asymptotic costate and the nonzero top rate}

Specialize to quasilinear preferences ($\mu = 1$) and assume the planner
places zero welfare weight on the top: $\psi(z) \to 0$ as $z\to\infty$ (a
purely redistributive planner).  For $z$ large enough that $\psi(z) = 0$:
\begin{equation}\label{eq:pareto_lambda}
  \lambda(z) = \int_z^\infty \bigl[\psi(\zeta)-1\bigr] f(\zeta)\,d\zeta
  \;\approx\; -\int_z^\infty f(\zeta)\,d\zeta
  = -(1-F(z))
  = -\!\left(\frac{\zlow}{z}\right)^\alpha < 0.
\end{equation}
The costate is \emph{negative} for high types and decays to zero at the Pareto
rate $z^{-\alpha}$.  Negative $\lambda(z)$ means the value function $V$ is
\emph{decreasing} in the inherited utility $\calU$: giving more utility to a
high-productivity type is socially costly because it tightens the incentive
constraint for all lower types, requiring the planner to give them more utility
to maintain IC.

\medskip
\begin{remark}[Negative costate and redistribution]
In the bounded-support problem, $\lambda(z) \geq 0$ everywhere under
redistributive weights (the costate equals the redistribution pressure from
above, which is non-negative).  With Pareto unbounded support and zero tail
weight, $\lambda(z) < 0$ for $z$ above the crossing point $\psi(z^{**}) = 1$.
The sign reversal signals that the planner is willing to earn revenue from top
types at the cost of distorting their labor supply---exactly what a positive
top marginal tax rate achieves.
\end{remark}

\subsection{The Diamond--Saez formula from the asymptotic HJB}

Substituting the asymptotic costate \eqref{eq:pareto_lambda} into the tax
formula derived from the HJB first-order condition (the general formula
\eqref{eq:M_taxformula} applied to the quasilinear case with constant
elasticity $e = 1/\gamma$):
\begin{align}
  \frac{T'}{1-T'}
  &= -\frac{1+\varepsilon^u}{\varepsilon^c} \cdot
     \frac{\lambda(z)}{z\,f(z)\,\mu} \notag \\
  &\approx \frac{1+\gamma}{\phantom{x}} \cdot
     \frac{1-F(z)}{z\,f(z)} \notag \\
  &= \frac{1+\gamma}{\alpha},
  \label{eq:pareto_wedge}
\end{align}
where the second line uses $\lambda(z)/\mu \approx -(1-F(z))$ and the third
uses the Pareto inverse hazard rate \eqref{eq:pareto_haz}.  Since
$(1+\varepsilon^u)/\varepsilon^c = (1+1/e)/(1/e) = 1+\gamma$ for quasilinear
preferences, and converting to the earnings Pareto
parameter via $\alpha = \alpha_y(1+e)/e = \alpha_y(1+\gamma)$:
\begin{equation}\label{eq:pareto_DS}
  \frac{T'_\infty}{1-T'_\infty}
  = \frac{1+\gamma}{\alpha_y(1+\gamma)}
  = \frac{1}{\alpha_y\, e},
\end{equation}
giving the Diamond--Saez formula:
\begin{equation}\label{eq:DS}
  \boxed{T'_\infty = \frac{1}{1+\alpha_y\, e}.}
\end{equation}

\subsection{Interpretation through the HJB lens}

Formula \eqref{eq:DS} can be read directly from three elements of the HJB:
\begin{enumerate}[(i)]
  \item \textbf{Costate decay rate $(\alpha)$.}  $-\lambda(z)/[1-F(z)] \to 1$:
    the costate decays at exactly the rate of the Pareto tail.  A heavier tail
    ($\alpha$ smaller) means the costate $\lambda(z)$ is numerically large
    relative to density, so the incentive-constraint distortion is large and
    the tax rate is high.

  \item \textbf{Inverse hazard rate $(1/\alpha)$.}  The Pareto inverse hazard
    rate \eqref{eq:pareto_haz} is the ratio $(1-F)/(zf)$ that converts the
    magnitude of $\lambda$ per person into the tax wedge per unit of income.
    With a thicker tail (smaller $\alpha$), more types are affected by any
    given marginal rate, so the revenue gain from a rate increase is large
    relative to the behavioral distortion at the margin.

  \item \textbf{Elasticity $(e)$.}  The factor $(1+\varepsilon^u)/\varepsilon^c
    = (1+\gamma)$ multiplies the inverse hazard rate.  A more elastic labor
    supply (smaller $1/\gamma$, i.e., larger $e$) amplifies the behavioral
    revenue loss from a rate increase, pushing the optimal rate down.
\end{enumerate}
The HJB perspective makes these three forces transparent: they are the
gradient of the value function ($\lambda$), the geometry of the type
distribution ($f$ and $F$), and the curvature of the flow payoff (the
elasticity terms in the maximand).

\subsection{Zero top rate under bounded support vs.\ positive top rate under
Pareto: a unified view}

In both cases the tax formula is
$T'/(1-T') = -(1+\varepsilon^u)/\varepsilon^c \cdot \lambda(z)/(zf(z)\mu)$.
The difference is entirely in the behavior of $\lambda(z)/[zf(z)]$:
\begin{itemize}
  \item \textbf{Bounded support:}  $1-F(z) \to 0$ as $z\to\zbar$ faster
    than density $f(z) > 0$, so $(1-F(z))/(z\,f(z))\to 0$.  Combined with
    $\lambda(\zbar)=0$, the tax formula gives $T'(\zbar) = 0$.

  \item \textbf{Pareto tail:}  $(1-F(z))/(z\,f(z)) \to 1/\alpha > 0$ and
    $\lambda(z)\approx -(1-F(z))$, so the ratio $\lambda(z)/(zf(z)) \to
    -1/\alpha \neq 0$.  The tax formula gives $T'_\infty = 1/(1+\alpha_y e)>0$.
\end{itemize}
The zero-rate result and the Diamond--Saez formula are thus two limiting cases
of the same HJB first-order condition, differing only in the asymptotic
behavior of the Pareto inverse hazard rate.

\subsection{Asymptotic value function under Pareto}

Under the linear guess the value function is $V(u,z;\mu) = \lambda(z)\,u +
K(z;\mu)$.  For large $z$:
\begin{equation}\label{eq:pareto_V}
  \lambda(z) \approx -(\zlow/z)^\alpha, \qquad
  K(z;\mu) \approx -\int_z^\infty \mathcal{F}(\zeta)\,d\zeta
\end{equation}
where $\mathcal{F}(\zeta) = \mu[y^*(\zeta)-C(\calU^*(\zeta),y^*(\zeta))]f(\zeta)$
is the flow resource surplus at $\zeta$.  Both $\lambda(z)$ and $K(z;\mu)$
decay polynomially at rate $z^{-\alpha}$, consistent with the transversality
condition $V(u,z;\mu) \to 0$.  The rate of decay is the Pareto shape
$\alpha$: heavier tails (smaller $\alpha$) slow the decay of the continuation
value, put more weight on incentive constraints at high types, and push up the
optimal tax rate.

\subsection{Table: bounded support vs.\ Pareto tail}

\begin{table}[h!]
\centering
\caption{Mirrlees HJB: Bounded vs.\ Pareto-tailed productivity}
\label{tab:pareto}
\renewcommand{\arraystretch}{1.45}
\begin{tabular}{@{} p{5.8cm} p{5.8cm} @{}}
\toprule
\textbf{Bounded support $\calZ=[\zlow,\zbar]$} &
\textbf{Pareto tail $\calZ=[\zlow,\infty)$} \\
\midrule
Upper BC: $\lambda(\zbar)=0$ (point condition) &
Upper BC: $\lambda(\infty)=0$ (transversality) \\
$(1-F(\zbar))/(z_{\bar{z}}\,f(\zbar)) = 0$ &
$(1-F(z))/(zf(z)) \to 1/\alpha > 0$ \\
$\lambda(z)/[zf(z)] \to 0$ as $z\to\zbar$ &
$\lambda(z)/[zf(z)] \to -1/\alpha < 0$ \\
Zero top rate: $T'(\zbar)=0$ &
Positive top rate: $T'_\infty=1/(1+\alpha_y e)>0$ \\
Value fn.\ decay: $V(u,z)\to 0$ at finite $\zbar$ &
Value fn.\ decay: $V(u,z)\sim z^{-\alpha}\to 0$ \\
Costate: $\lambda(z)\geq 0$ is hump-shaped &
Costate: $\lambda(z)$ turns negative for large $z$ \\
\bottomrule
\end{tabular}
\end{table}

\begin{remark}[Empirical calibration]
The Diamond--Saez formula \eqref{eq:DS} requires two empirical inputs:
$\alpha_y$ (the Pareto shape of the earnings distribution) and $e$ (the
compensated elasticity of taxable income).  Saez (2001) surveys estimates
suggesting $\alpha_y \approx 2$ for the United States, and a benchmark value
of $e \approx 0.25$ for broad income, giving $T'_\infty \approx 1/(1+2\times
0.25) = 2/3 \approx 67\%$.  With higher elasticity, say $e = 0.5$, the
formula gives $T'_\infty = 1/(1+1) = 50\%$.  Diamond and Saez (2011)
conclude that $T'_\infty$ in the range $50$--$70\%$ is consistent with the
empirical evidence, well above the top marginal rates prevailing in most
countries at the time.
\end{remark}

%% ===========================================================================
\appendix
\section{The Hamiltonian and HJB for Mirrlees: Equivalence}
\label{app:hamilton_hjb}
%% ===========================================================================

The Hamiltonian approach of the lecture notes defines
\begin{equation}\label{eq:M_Hamilton}
  \mathcal{H}(\calU, y, \lambda;\, z) =
  \bigl[\psi(z)\,\calU + \mu(y - C(\calU, y))\bigr] f(z)
  - \lambda\, U_n\!\bigl(C(\calU, y), y/z\bigr)\frac{y}{z^2}.
\end{equation}
The HJB \eqref{eq:M_HJBcompact} is the Hamilton--Jacobi equation
\begin{equation}\label{eq:M_HJ}
  -V_z(u, z) = \mathcal{H}\bigl(u,\, y^*(u,z),\, V_u(u,z);\, z\bigr),
\end{equation}
where $y^*(u,z) = \argmax_y \mathcal{H}(u, y, V_u(u,z); z)$.  Along the
optimal path, the costate $\lambda(z) = V_u(\calU^*(z), z)$ satisfies
$d\lambda/dz = -\mathcal{H}_\calU$ (equation \eqref{eq:M_costate_general}
above), confirming that the Hamiltonian and HJB approaches are equivalent.

\end{document}
